% Obsahový zdroj: import_hodiny/29_Uvod_do_kombinatoriky.tex a odpovídající PDF (obě strany).
% P1–P4: zdrojové Příklady 1–4; C1–C5: zdrojové Procvičování 1–5.
% Matematický obsah zachován; výklad, odvození a řešení pedagogicky rozvedeny.
% Upřesnění a opravy: viz KMB_29-41_Kontrola.md; původní import zůstává beze změny.
% Rozpis: 4_rocnik/tematicke_plany/Rozpis_vyuky_MAT_4_rocnik.tex, hodina 29.
% Plán: 4_rocnik/tematicke_plany/SPS.pdf, oddíly 2.1–2.3.
% Sazba: sablony/pripravy/TEMPLATE_Priprava.tex a společný priprava.sty.
% Vlastní doplnění: didaktický výklad, zdůvodnění, metodika; žádné nové číslované úlohy.
% Učebnice/09_Kombinatorika nemá obsahovou kapitolu; místní kombinatorické listy ani Realisticky.cz nejsou dostupné.
\documentclass[11pt,a4paper]{article}
% Spolecny styl pro ucitelske pripravy MAT (Work/MAT)
% Verze 1.1, 2026-10-23
\NeedsTeXFormat{LaTeX2e}
\ProvidesPackage{priprava}[2026/10/23 v1.1 Sablona ucitelskych priprav MAT]

% Zakladni typografie a cestina
\RequirePackage[T1]{fontenc}
\RequirePackage[utf8]{inputenc}
\RequirePackage[czech]{babel}
\RequirePackage{lmodern}
\RequirePackage{microtype}
\RequirePackage{amsmath,amssymb}
\RequirePackage{xcolor}
\RequirePackage{enumitem}
\RequirePackage[most]{tcolorbox}
\RequirePackage{tabularx,booktabs,array}
\RequirePackage{fancyhdr}
\RequirePackage{needspace}
\RequirePackage{geometry}
\RequirePackage[hidelinks]{hyperref}

\geometry{a4paper,top=17mm,bottom=18mm,left=23mm,right=23mm}
\definecolor{notegray}{gray}{0.39}
\setlength{\parindent}{0pt}
\setlength{\parskip}{2.5pt plus 1pt minus 0.5pt}
\setlist[itemize]{topsep=2pt,itemsep=1pt,parsep=0pt,leftmargin=1.45em}
\setlist[enumerate]{topsep=2pt,itemsep=2pt,parsep=0pt,leftmargin=1.65em}
\setcounter{secnumdepth}{0}
\renewcommand{\arraystretch}{1.12}

% Zahlavi: vlevo rocnik | globalni tema, vpravo cislo: nazev hodiny
% Zapati zustava zcela prazdne.
\pagestyle{fancy}
\fancyhf{}
\renewcommand{\headrulewidth}{0.35pt}
\renewcommand{\footrulewidth}{0pt}
\setlength{\headheight}{14pt}
\newcommand{\nastavPripravu}[4]{%
  \fancyhead[L]{\small #1\enspace\textbar\enspace #2}%
  \fancyhead[R]{\small #3: #4}%
}

% Minimalisticke cernobile boxy. Kazda definice zvlast.
\tcbset{commonlesson/.style={enhanced,breakable=false,colback=white,colframe=black,
  colbacktitle=white,coltitle=black,boxrule=0.45pt,
  arc=0.8mm,outer arc=0.8mm,left=3mm,right=3mm,top=1.6mm,bottom=1.6mm,
  toptitle=1mm,bottomtitle=0.7mm,fonttitle=\bfseries,
  attach boxed title to top left={xshift=2.5mm,yshift=-1.15mm},
  boxed title style={colback=white,colframe=white,boxrule=0pt,left=1mm,right=1mm}}}
\newtcolorbox{zadaniBox}[1][]{commonlesson,title={Zadání},#1}
\newtcolorbox{prikladBox}[1][]{commonlesson,title={Příklad},#1}
\newtcolorbox{definiceBox}[1][]{commonlesson,title={Definice},#1}

% Spolecny box pro zadani a jeho reseni, bez mezery mezi castmi.
% Pouziti:
%   \begin{prikladReseniBox}
%     Zadani ulohy.
%     \tcblower
%     \begin{align*} ... \end{align*}
%   \end{prikladReseniBox}
% Spodni cast ma jemny sedy podklad a delici linku; lze v ni sazet
% vice radku matematiky. Obe casti maji jen jeden spolecny ramecek.
\tcbset{solutionlesson/.style={
  commonlesson,
  enhanced jigsaw,
  breakable,
  colbacklower=black!3,
  segmentation style={solid,draw=black!25,line width=0.35pt},
  fontlower=\small\color{notegray},
  before lower={\noindent\textbf{Řešení:}\enspace},
  leftlower=3mm,rightlower=3mm
}}
\newtcolorbox{prikladReseniBox}[1][]{solutionlesson,title={Příklad},#1}
\newtcolorbox{zadaniReseniBox}[1][]{solutionlesson,title={Zadání},#1}

% Pomucky pro autora pripravy: sedive poznamky; \reseni pro starsi pripravy
\newcommand{\poznamka}[1]{\par\begingroup\small\color{notegray}\itshape
  \textnormal{\textbf{Poznámka pro učitele.}}\enspace #1\par\endgroup}
\newcommand{\reseni}[1]{\par\begingroup\small\color{notegray}
  \textbf{Očekávané řešení:}\enspace #1\par\endgroup}
\newcommand{\blok}[1]{\needspace{5\baselineskip}\subsection*{#1}}
\newcommand{\vyrok}[1]{\uv{#1}}

\nastavPripravu{4. ročník}{Kombinatorika}{29}{Úvod do kombinatoriky}
\begin{document}


Začni dvěma situacemi se stejnými kartičkami: seřadit je do řady a vybrat z nich dvojici. Nech žáky nejprve vyslovit odhad a potom vypsat možnosti. Rozdíl mezi výsledky musí vysvětlit významem pořadí.
\poznamka{Předpoklady: konečná množina, dvojice, jednoduché násobení. Cíl: žák přesně popíše jeden výsledek, rozliší pořadí a opakování a provede úplný výpis bez duplicit. Pojmy permutace, variace a kombinace zatím nevyžaduj jako názvy vzorců.}

\blok{Výklad a odvození}
\begin{definiceBox}
\textbf{Kombinatorika} zkoumá mimo jiné počty výběrů a uspořádání prvků podle daných podmínek. Před počítáním musí být určeno, kdy dva výsledky považujeme za stejné.
\end{definiceBox}
\begin{definiceBox}
\textbf{Uspořádaný výběr} rozlišuje pozice prvků: dvojice $(A,B)$ a $(B,A)$ jsou pro různé prvky $A,B$ různé výsledky.
\end{definiceBox}
\begin{definiceBox}
\textbf{Neuspořádaný výběr} rozlišuje pouze vybrané prvky, nikoli jejich pořadí: $\{A,B\}=\{B,A\}$.
\end{definiceBox}
\begin{definiceBox}
\textbf{Výběr s opakováním} dovoluje použít tentýž prvek vícekrát.
\end{definiceBox}
\begin{definiceBox}
Při \textbf{výběru bez opakování} se každý prvek vyskytne nejvýše jednou.
\end{definiceBox}

Výpis organizuj podle prvního prvku, potom podle druhého. Každá možnost musí mít právě jedno místo v seznamu. Úplnost prokážeme tím, že jsme u každého začátku prošli všechna dovolená pokračování; odlišnost tím, že dvě větve dávají různé výsledky.

Ptej se postupně: Co tvoří jeden výsledek? Změní záměna pořadí výsledek? Lze tentýž prvek použít znovu? Použijeme všechny prvky, nebo jen část? Například komise je skupina osob, zatímco první a druhé místo jsou dvě odlišné role.


\poznamka{Průběh: 0–8 min kartičky a P1; 8–20 min P2–P3 a rozdíl výběr/uspořádání; 20–30 min P4 a význam opakování; 30–41 min C1–C5; 41–45 min kontrola a slovní shrnutí. U P4 je násobení zatím zkrácený zápis větvení; obecné pravidlo přijde v hodině 30.}

\blok{Společně řešené příklady}

% Zdroj P1: import_hodiny/29_Uvod_do_kombinatoriky.tex, Příklad 1.
\begin{prikladReseniBox}[title={Příklad P1}]
Tři označené kartičky \(A,B,C\) postavíme do řady. Kolik vznikne pořadí?
\tcblower
Rozlišíme první, druhé a třetí místo. Pro každý první znak vypíšeme obě pořadí zbylých kartiček:
\[
 ABC,\ ACB;\qquad BAC,\ BCA;\qquad CAB,\ CBA.
\]
Začátek je vždy $A$, $B$ nebo $C$, a proto žádná možnost nechybí. V seznamu není žádné pořadí dvakrát. Celkem je $6$ pořadí.
\end{prikladReseniBox}


% Zdroj P2: import_hodiny/29_Uvod_do_kombinatoriky.tex, Příklad 2.
\begin{prikladReseniBox}[title={Příklad P2}]
Ze stejných tří kartiček vybereme dvě bez rozlišení pořadí.
\tcblower
Přípustné dvojice jsou \(\{A,B\},\{A,C\},\{B,C\}\). Nejsou to šestice: \(AB\) a \(BA\) označují stejný výběr. Výsledek: \(3\).
\end{prikladReseniBox}


% Zdroj P3: import_hodiny/29_Uvod_do_kombinatoriky.tex, Příklad 3.
\begin{prikladReseniBox}[title={Příklad P3}]
Čtyři různí žáci: vyber dva členy komise a zvlášť přiděl první a druhé místo.
\tcblower
Komise: \(\{A,B\},\{A,C\},\{A,D\},\{B,C\},\{B,D\},\{C,D\}\), tedy \(6\) možností. Při přidělení prvního a druhého místa má každá dvojice dvě pořadí, tedy \(12\) možností.
\end{prikladReseniBox}


% Zdroj P4: import_hodiny/29_Uvod_do_kombinatoriky.tex, Příklad 4.
\begin{prikladReseniBox}[title={Příklad P4}]
Tříznakový kód ze symbolů \(0,1,2\), s opakováním. Jaký počet očekáváš?
\tcblower
Pro začátek $0$ vypíšeme všechny dvojice dalších znaků: $00,01,02,10,11,12,20,21,22$. Je jich devět. Stejných devět pokračování má začátek $1$ i $2$, tedy $9+9+9=3\cdot3\cdot3=27$ kódů. Opakování je dovoleno; $000$ i $012$ jsou platné kódy. Každý kód jsme započítali právě v jedné skupině podle jeho začátku.
\end{prikladReseniBox}


\newpage
\blok{Samostatná práce a zásoba k procvičení}

Nech nejprve samostatně zapsat vlastní odpověď a její zdůvodnění. U C4–C5 se ptáme pouze na vlastnosti výběru, nikoli na jeho počet. Řešení ve spodních částech slouží učiteli.

% Zdroj C1: import_hodiny/29_Uvod_do_kombinatoriky.tex, Procvičování, položka 1.
\begin{zadaniReseniBox}[title={Úloha C1}]
Kolik různých dvojic vznikne výběrem ze čtyř lidí?
\tcblower
Každou neuspořádanou dvojici vypíšu jednou: \(AB,AC,AD,BC,BD,CD\). Celkem \(6\).
\end{zadaniReseniBox}

% Zdroj C2: import_hodiny/29_Uvod_do_kombinatoriky.tex, Procvičování, položka 2.
\begin{zadaniReseniBox}[title={Úloha C2}]
Kolika způsoby obsadí tři lidé tři židle v řadě?
\tcblower
První židle: \(3\) osoby, druhá \(2\), poslední \(1\), tedy \(3\cdot2\cdot1=6\).
\end{zadaniReseniBox}

% Zdroj C3: import_hodiny/29_Uvod_do_kombinatoriky.tex, Procvičování, položka 3.
\begin{zadaniReseniBox}[title={Úloha C3}]
Kolik výsledků mají dva po sobě jdoucí hody mincí?
\tcblower
Označme strany mince $L$ (líc) a $R$ (rub). Jeden výsledek je uspořádaná dvojice hodů: $LL,LR,RL,RR$; jsou čtyři.
\end{zadaniReseniBox}

% Zdroj C4: import_hodiny/29_Uvod_do_kombinatoriky.tex, Procvičování, položka 4.
\begin{zadaniReseniBox}[title={Úloha C4}]
U výběru tříčlenné komise ze sedmi žáků rozhodni: záleží na pořadí, lze opakovat žáka?
\tcblower
Nezáleží na pořadí; každý žák může v komisi figurovat nejvýše jednou. Jde o neuspořádaný výběr bez opakování.
\end{zadaniReseniBox}

% Zdroj C5: import_hodiny/29_Uvod_do_kombinatoriky.tex, Procvičování, položka 5.
\begin{zadaniReseniBox}[title={Úloha C5}]
U čtyřmístného PINu rozhodni: pořadí? opakování? počáteční nula?
\tcblower
Na pořadí záleží, opakování je dovoleno a počáteční nula ano: \(0123\) je platný PIN.
\end{zadaniReseniBox}


\poznamka{Typické chyby: dvojice $AB$ a $BA$ započtené dvakrát; opomenuté větve; záměna PINu a čísla. Zeptej se, proč kód $000$ není v P4 zakázán. Na závěr vyžaduj vlastními slovy rozdíl mezi komisí a pořadím na stupních vítězů. Hodina 30 zobecní výpis pravidlem součtu a součinu.}
\end{document}
